\documentclass[12pt,a4paper,oneside]{book}
\usepackage{amsmath}
\usepackage{amssymb}
\usepackage{amsthm}
\usepackage{graphicx}
\usepackage{latexsym}
\usepackage{color}
\usepackage{fancyhdr}
\usepackage{cite}
\usepackage{indentfirst}
\usepackage{tocloft}
\usepackage{multirow}
\usepackage{fontspec}
\usepackage{listings}
\usepackage{xcolor}

% Times New Roman
\setromanfont[
BoldFont=Times New Roman Bold.ttf,
ItalicFont=Times New Roman Italic.ttf,
BoldItalicFont=Times New Roman Bold Italic.ttf,
]{Times New Roman.ttf}

% Centering Chapter
\usepackage{sectsty}   
\chapterfont{\centering}

\pagestyle{myheadings}
\setlength{\parindent}{0.8in}

\usepackage[top=1.5in,bottom=1in,left=1.5in,right=1in]{geometry}

% Theorem style
\theoremstyle{plain}
\newtheorem{thm}{Theorem}[chapter]
\newtheorem{lem}[thm]{Lemma}
\newtheorem{cor}[thm]{Corollary}
\newtheorem{prop}[thm]{Proposition}
\newtheorem{rem}[thm]{Remark}
\newtheorem{ex}[thm]{Example}
\newtheorem{de}[thm]{Definition}
\renewcommand{\proof}{\textbf{Proof.}}

\numberwithin{equation}{chapter}
\DeclareMathOperator{\Var}{Var}
\DeclareMathOperator{\Ima}{Im}

\usepackage{float}

% Graph
\usepackage{graphics,graphicx}
\usepackage{calc}
\usepackage{tikz}
\usetikzlibrary{decorations.markings}
\tikzstyle{vertex}=[circle, draw, inner sep=2pt, minimum size=4pt]
\newcommand{\vertex}{\node[vertex]}
\newcounter{Angle}

% Line space
\renewcommand{\baselinestretch}{1.5}

\newcommand*{\QEDA}{\hfill\ensuremath{\large{\lozenge}}}
\newcommand*{\QEDAa}{\hfill\ensuremath{\square}}

\renewcommand{\bibname}{\centerline{\Large TÀI LIỆU THAM KHẢO}}

% Centering table of contents and list of table
\usepackage{tocloft}
\renewcommand{\contentsname}{\hfill\bfseries\Large MỤC LỤC \hfill}   
\renewcommand{\cftaftertoctitle}{\hfill}

\renewcommand{\listtablename}{\hfill\bfseries\Large DANH MỤC BẢNG} 
\renewcommand{\cftafterlottitle}{\hfill}

\renewcommand{\listfigurename}{\hfill\bfseries\Large DANH MỤC HÌNH ẢNH} 
\renewcommand{\cftafterloftitle}{\hfill}

\lstdefinestyle{sqlstyle}{
    language=SQL,
    basicstyle=\ttfamily\small,
    keywordstyle=\color{blue}\bfseries,
    stringstyle=\color{red},
    commentstyle=\color{gray}\itshape,
    numbers=left,
    numberstyle=\tiny\color{gray},
    stepnumber=1,
    breaklines=true,
    frame=single,
    backgroundcolor=\color{gray!10},
    tabsize=4,
}

\lstdefinestyle{csharpstyle}{
    language=[Sharp]C,
    basicstyle=\ttfamily\small,
    keywordstyle=\color{blue}\bfseries,
    stringstyle=\color{red!70!black},
    commentstyle=\color{green!50!black}\itshape,
    numbers=left,
    numberstyle=\tiny\color{gray},
    stepnumber=1,
    breaklines=true,
    frame=single,
    backgroundcolor=\color{gray!10},
    tabsize=4,
}

\begin{document}
\thispagestyle{empty}

\begin{figure}[h!]
\begin{center}
\includegraphics[width=7.5cm]{HCMCOU_Logo.png}
\end{center}
\end{figure}

\begin{center}
\large{\textbf{HỆ THỐNG QUẢN LÝ RẠP PHIM}}
\end{center}

\vskip0.3cm

\begin{center}
\textbf{Thành viên}
\end{center}

\vskip0.6cm

\begin{center}
\textbf{Trần Thanh Tung}\\
\textbf{Họ Đệm Tên 2}\\
\textbf{Họ Đệm Tên 3}
\end{center}

\vskip3.5cm

\begin{center}
\textbf{BÁO CÁO ĐỒ ÁN MÔN HỌC}
\end{center}

\begin{center}
\textbf{LẬP TRÌNH CƠ SỞ DỮ LIỆU}
\end{center}

\vskip2cm

\begin{center}
\textbf{TRƯỜNG ĐẠI HỌC MỞ TP.HCM (HCMCOU)}
\end{center}

\newpage
\pagenumbering{roman}

\begin{table}[h]
    \begin{tabular}{ll}
        Tên đồ án       & Hệ thống Quản lý Rạp Phim (Cinema Management System) \\
        Thành viên      & Trần Thanh Tung \\
        Môn học:        & Lập trình Cơ sở dữ liệu \\
        Người hướng dẫn & ThS. Phạm Hoàng An \\
        Năm học         & 2025-2026 \\
    \end{tabular}
\end{table}

\begin{center}
  \large{\textbf{GIỚI THIỆU}}\\
\end{center}
\addcontentsline{toc}{chapter}{GIỚI THIỆU}

Đồ án này trình bày quá trình thiết kế và xây dựng hệ thống Quản lý Rạp Phim --- một ứng dụng web cho phép quản lý phim, suất chiếu, đặt vé trực tuyến và thanh toán. Hệ thống được xây dựng trên nền tảng \textbf{ASP.NET Core MVC (.NET 8)} với kiến trúc \textbf{3 lớp} (Presentation -- Business Logic -- Data Access), sử dụng \textbf{Entity Framework Core} kết hợp \textbf{ADO.NET} và cơ sở dữ liệu \textbf{SQL Server (T-SQL)}.

Người dùng cuối có thể duyệt danh sách phim, xem chi tiết, chọn suất chiếu, chọn ghế ngồi, thêm dịch vụ bắp nước/combo, xác nhận thanh toán và nhận vé điện tử (E-Ticket) kèm mã QR. Phía quản trị viên có thể quản lý toàn bộ danh mục phim, suất chiếu, dịch vụ và xem báo cáo doanh thu theo phim thông qua biểu đồ Chart.js tích hợp.

Ngoài ra, hệ thống cũng phát triển \textbf{RESTful Web API} với đầy đủ các endpoint CRUD, hỗ trợ export/import dữ liệu qua định dạng XML sử dụng \textbf{LINQ to XML}.

Báo cáo được tổ chức theo các chương: Chương 1 trình bày tổng quan hệ thống; Chương 2 phân tích và thiết kế cơ sở dữ liệu T-SQL; Chương 3 mô tả kiến trúc phần mềm 3 lớp; Chương 4 trình bày các kỹ thuật triển khai bao gồm Entity Framework, LINQ, ADO.NET (DataSet/DataAdapter), LINQ to XML và Web API; Chương 5 kết luận và hướng phát triển.

\noindent{\textbf{Từ khóa:}} ASP.NET Core MVC, Entity Framework Core, SQL Server T-SQL, ADO.NET, LINQ, LINQ to XML, RESTful Web API, Kiến trúc 3 lớp, Quản lý rạp phim.

\newpage
\tableofcontents

\newpage
\addcontentsline{toc}{chapter}{DANH MỤC HÌNH ẢNH}
\listoffigures

\newpage
\addcontentsline{toc}{chapter}{DANH MỤC BẢNG}
\listoftables

\newpage
\pagenumbering{arabic}

\chapter{Giới thiệu chung}

\noindent Bảng~\ref{table1} liệt kê danh sách các thành viên tham gia thực hiện đồ án và phân công nhiệm vụ.

\begin{table}[H]
    \begin{center}
        \begin{tabular}{|c|l|l|c|}
            \hline
            \textbf{MSSV} & \textbf{Họ tên} & \textbf{Nhiệm vụ} & \textbf{Ghi chú} \\
            \hline
            2351010232 & Trần Thanh Tung & Backend -- BUS \& DAL Layer, API & \\
            \hline
            & Họ Đệm Tên 2 & Thiết kế CSDL, T-SQL & \\
            \hline
            & Họ Đệm Tên 3 & Frontend -- Razor Views, Admin UI & \\
            \hline
        \end{tabular}
    \end{center}
    \caption{Danh sách thành viên thực hiện đồ án}\label{table1}
\end{table}

\vspace{1cm}

\noindent Hệ thống Quản lý Rạp Phim cung cấp hai nhóm chức năng chính:

\begin{itemize}
    \item \textbf{Khách hàng:} Đăng ký, đăng nhập, duyệt danh sách phim, chọn suất chiếu, chọn ghế, thêm dịch vụ bắp nước/combo, thanh toán và nhận E-Ticket có mã QR. Sau khi đăng nhập, người dùng có thể cập nhật thông tin cá nhân (Profile) và đổi mật khẩu.
    \item \textbf{Quản trị viên:} Quản lý phim (thêm, sửa, xóa, upload ảnh), quản lý suất chiếu (tự động tính giờ kết thúc), quản lý dịch vụ bắp nước (CRUD), xem báo cáo doanh thu theo từng bộ phim qua biểu đồ.
\end{itemize}

\noindent Hình~\ref{fig1} là logo của trường, được sử dụng trong giao diện ứng dụng.

\begin{figure}[H]
    \includegraphics[width=3.5cm]{HCMCOU_Logo.png}
    \centering
    \caption{Logo Trường Đại học Mở TP.HCM}\label{fig1}
\end{figure}

\noindent Công nghệ và công cụ sử dụng trong dự án được tóm tắt trong Bảng~\ref{table_tech}.

\begin{table}[H]
    \begin{center}
        \begin{tabular}{|l|l|}
            \hline
            \textbf{Thành phần} & \textbf{Công nghệ / Phiên bản} \\
            \hline
            Ngôn ngữ lập trình & C\# (.NET 8.0) \\
            \hline
            Web Framework      & ASP.NET Core MVC 8.0 \\
            \hline
            ORM                & Entity Framework Core 8.0.0 \\
            \hline
            Truy cập dữ liệu   & ADO.NET (SqlDataAdapter, DataSet, DataTable) \\
            \hline
            Cơ sở dữ liệu     & Microsoft SQL Server Express 2019+ \\
            \hline
            Frontend           & Razor View, Bootstrap 5.3.2, jQuery 3.7.1 \\
            \hline
            Biểu đồ            & Chart.js (latest) \\
            \hline
            Hộp thoại xác nhận & SweetAlert2 v11 \\
            \hline
            IDE                & Visual Studio 2022 (v17.x) \\
            \hline
            Quản lý phiên bản  & Git / GitHub \\
            \hline
        \end{tabular}
    \end{center}
    \caption{Bảng công nghệ sử dụng}\label{table_tech}
\end{table}

\chapter{Phân tích và Thiết kế Cơ sở Dữ liệu (T-SQL)}

\section{Sơ đồ thực thể mối quan hệ (ERD)}\label{Sec2.1}

\begin{de}\label{def1}
    \textbf{Các thực thể chính} trong hệ thống bao gồm: \texttt{NguoiDung}, \texttt{Phim}, \texttt{LoaiPhim}, \texttt{Phong}, \texttt{SuatChieu}, \texttt{Ghe}, \texttt{HoaDon}, \texttt{Ve}, \texttt{DichVu} và \texttt{ChiTietDV}. Mỗi thực thể ánh xạ tương ứng với một bảng trong cơ sở dữ liệu \texttt{QuanLyRapPhim}. Toàn bộ ánh xạ này được thể hiện trong \texttt{QuanLyRapPhimContext.cs} thông qua phương thức \texttt{OnModelCreating()}.
\end{de}

\vspace{0.5cm}

\noindent Cơ sở dữ liệu gồm 10 bảng chính với các mối quan hệ khóa ngoại được mô tả trong Bảng~\ref{table_db}.

\begin{table}[H]
    \begin{center}
        \begin{tabular}{|l|l|l|}
            \hline
            \textbf{Bảng (Table)} & \textbf{Khóa chính (PK)} & \textbf{Mô tả} \\
            \hline
            NguoiDung   & MaND (MaNd)  & Tài khoản người dùng, IsAdmin phân quyền \\
            \hline
            Phim        & MaPhim       & Thông tin bộ phim + thể loại \\
            \hline
            LoaiPhim    & MaLoai       & Thể loại phim (Hành động, Hoạt hình...) \\
            \hline
            Phong       & MaPhong      & Phòng chiếu với số lượng ghế \\
            \hline
            SuatChieu   & MaSuat       & Suất chiếu: ngày, giờ, giá vé, trạng thái \\
            \hline
            Ghe         & MaGhe        & Ghế trong phòng, loại ghế \\
            \hline
            HoaDon      & MaHD (MaHd)  & Hóa đơn đặt vé, liên kết NguoiDung \\
            \hline
            Ve          & MaVe         & Vé xem phim, khóa duy nhất (MaSuat, MaGhe) \\
            \hline
            DichVu      & MaDV (MaDv)  & Dịch vụ bắp nước/combo, số lượng tồn \\
            \hline
            ChiTietDV   & (MaHD, MaDV) & Chi tiết dịch vụ trong hóa đơn, giá bán lúc mua \\
            \hline
        \end{tabular}
    \end{center}
    \caption{Danh sách các bảng trong cơ sở dữ liệu \texttt{QuanLyRapPhim}}\label{table_db}
\end{table}

\section{Các ràng buộc và đối tượng lập trình T-SQL}\label{Sec2.2}

\begin{lem}\label{lem1}
    Các \textbf{ràng buộc toàn vẹn dữ liệu} được áp dụng trong \texttt{QuanLyRapPhimContext.cs} qua \texttt{OnModelCreating()} bao gồm:
    \begin{itemize}
        \item \textbf{Khóa ngoại (FK):} \texttt{Ve.MaGhe} $\to$ \texttt{Ghe.MaGhe}; \texttt{Ve.MaHD} $\to$ \texttt{HoaDon.MaHD}; \texttt{Ve.MaSuat} $\to$ \texttt{SuatChieu.MaSuat}; \texttt{SuatChieu.MaPhim} $\to$ \texttt{Phim.MaPhim}; \texttt{Phim.MaLoaiPhim} $\to$ \texttt{LoaiPhim.MaLoai}; v.v.
        \item \textbf{Duy nhất (Unique Index):} \texttt{NguoiDung.TaiKhoan} là duy nhất (\texttt{UQ\_\_NguoiDun\_...}); cặp \texttt{(Ve.MaSuat, Ve.MaGhe)} không trùng lặp (\texttt{UC\_SuatChieu\_Ghe}).
        \item \textbf{Giá trị mặc định (Default):} \texttt{HoaDon.NgayLap = GETDATE()}; \texttt{HoaDon.TongTien = 0}; \texttt{Ghe.DaDat = false}; \texttt{SuatChieu.TrangThai = N'Sắp chiếu'}.
        \item \textbf{Kiểu tiền tệ:} Các cột \texttt{DonGia}, \texttt{GiaVe}, \texttt{TongTien}, \texttt{DonGiaBan} sử dụng kiểu \texttt{money} của SQL Server.
    \end{itemize}
\end{lem}

\vspace{0.5cm}

\subsection{View T-SQL}

\noindent Hệ thống định nghĩa View \texttt{vw\_DoanhThuTheoPhim} để thống kê doanh thu tổng hợp theo từng bộ phim, được dùng tại Admin \textit{Báo cáo} và API endpoint \texttt{/api/phimapi/adonet-dataset}:

\begin{lstlisting}[style=sqlstyle, caption={View thống kê doanh thu theo phim (\texttt{vw\_DoanhThuTheoPhim})}]
CREATE VIEW vw_DoanhThuTheoPhim AS
SELECT 
    p.MaPhim,
    p.TenPhim,
    COUNT(DISTINCT hd.MaHD) AS SoLuongHoaDon,
    ISNULL(SUM(hd.TongTien), 0) AS TongDoanhThu
FROM Phim p
LEFT JOIN SuatChieu s  ON p.MaPhim = s.MaPhim
LEFT JOIN Ve v         ON s.MaSuat = v.MaSuat
LEFT JOIN HoaDon hd    ON v.MaHD = hd.MaHD
GROUP BY p.MaPhim, p.TenPhim;
\end{lstlisting}

\subsection{Function T-SQL}

\noindent Function \texttt{fn\_TinhTongTienHoaDon} tính tổng tiền của một hóa đơn bao gồm cả tiền vé và dịch vụ đi kèm:

\begin{lstlisting}[style=sqlstyle, caption={Scalar Function tính tổng tiền hóa đơn}]
CREATE FUNCTION fn_TinhTongTienHoaDon(@MaHD INT)
RETURNS MONEY
AS
BEGIN
    DECLARE @TongTienVe MONEY = 0;
    DECLARE @TongTienDV MONEY = 0;
    
    SELECT @TongTienVe = ISNULL(SUM(GiaVe), 0)
    FROM Ve WHERE MaHD = @MaHD;
    
    SELECT @TongTienDV = ISNULL(SUM(DonGiaBan * SoLuong), 0)
    FROM ChiTietDV WHERE MaHD = @MaHD;
    
    RETURN @TongTienVe + @TongTienDV;
END
\end{lstlisting}

\subsection{Stored Procedure T-SQL}

\noindent Stored Procedure \texttt{sp\_LayDanhSachPhimDangChieu} trả về các phim có suất chiếu từ hôm nay trở đi, được gọi từ \texttt{PhimBUS.LayDanhSachPhimDangChieu\_SP()} và endpoint \texttt{GET /api/phimapi} qua \texttt{FromSqlRaw}:

\begin{lstlisting}[style=sqlstyle, caption={Stored Procedure lấy phim đang chiếu}]
CREATE PROCEDURE sp_LayDanhSachPhimDangChieu
AS
BEGIN
    SET NOCOUNT ON;
    SELECT p.MaPhim, p.TenPhim, p.ThoiLuong, 
           p.MaLoaiPhim, p.GioiHanTuoi, 
           p.NgayKhoiChieu, p.Hinh
    FROM Phim p
    WHERE EXISTS (
        SELECT 1 FROM SuatChieu s 
        WHERE s.MaPhim = p.MaPhim 
          AND CAST(s.NgayChieu AS DATE) >= CAST(GETDATE() AS DATE)
    );
END
\end{lstlisting}

\subsection{Trigger T-SQL}

\noindent Trigger \texttt{trg\_NganXoaHoaDonCoVe} sử dụng \texttt{INSTEAD OF DELETE} để bảo vệ các hóa đơn đã được xuất vé:

\begin{lstlisting}[style=sqlstyle, caption={Trigger ngăn xóa hóa đơn đã có vé}]
CREATE TRIGGER trg_NganXoaHoaDonCoVe
ON HoaDon
INSTEAD OF DELETE
AS
BEGIN
    IF EXISTS (
        SELECT 1 FROM Ve v
        INNER JOIN deleted d ON v.MaHD = d.MaHD
    )
    BEGIN
        RAISERROR(N'Khong the xoa hoa don da co du lieu xuat ve!',
                  16, 1);
        ROLLBACK TRANSACTION;
        RETURN;
    END
    DELETE HD FROM HoaDon HD
    INNER JOIN deleted d ON HD.MaHD = d.MaHD;
END
\end{lstlisting}

\subsection{Transaction (Xử lý giao tác)}

\begin{thm}\label{thm_transaction}
    Toàn bộ tiến trình thanh toán (lưu \texttt{HoaDon}, \texttt{Ve} và \texttt{ChiTietDV}) được bao bọc trong một \textbf{Database Transaction} của Entity Framework Core. Nếu bất kỳ bước nào thất bại, toàn bộ sẽ bị \texttt{Rollback}, đảm bảo tính nguyên tử (Atomicity).
\end{thm}

\begin{lstlisting}[style=csharpstyle, caption={Transaction trong \texttt{HoaDonBUS.LuuVaThanhToan()}}]
public int LuuVaThanhToan(CartItemDTO cart, string username)
{
    using (var transaction = _context.Database.BeginTransaction())
    {
        try
        {
            // 1. Tạo Hóa đơn
            var hoaDon = new HoaDon {
                MaNd = user.MaNd,
                NgayLap = DateTime.Now,
                TongTien = cart.TongTien
            };
            _context.HoaDons.Add(hoaDon);
            _context.SaveChanges();
            
            // 2. Lưu Vé cho từng ghế đã chọn
            foreach (var ghe in dsGheTrongPhong)
            {
                _context.Ves.Add(new Ve {
                    MaHd = hoaDon.MaHd,
                    MaSuat = cart.MaSuat,
                    MaGhe = ghe.MaGhe,
                    GiaVe = suatChieu.GiaVe
                });
            }
            
            // 3. Lưu Chi tiết Dịch vụ (Combo bắp nước)
            foreach (var dv in cart.DichVus)
            {
                _context.ChiTietDvs.Add(new ChiTietDv {
                    MaHd = hoaDon.MaHd,
                    MaDv = dv.MaDV,
                    SoLuong = dv.SoLuong,
                    DonGiaBan = dv.DonGia
                });
            }
            _context.SaveChanges();
            transaction.Commit();   
            return hoaDon.MaHd;
        }
        catch (Exception)
        {
            transaction.Rollback(); 
            return 0;
        }
    }
}
\end{lstlisting}

\chapter{Kiến trúc Phần mềm (3-Layer Architecture)}

\noindent Hệ thống được tổ chức theo kiến trúc \textbf{3 lớp} tách biệt hoàn toàn gồm bốn dự án (project) riêng trong một Solution:

\begin{itemize}
    \item \textbf{Cinema.DAL} --- Lớp truy cập dữ liệu (EF Core DbContext, Models, ADO.NET)
    \item \textbf{Cinema.BUS} --- Lớp nghiệp vụ (Business Logic, Interface + Implementation)
    \item \textbf{Cinema.DTO} --- Đối tượng trao đổi dữ liệu giữa các lớp
    \item \textbf{Cinema.Web} --- Lớp giao diện (ASP.NET Core MVC + RESTful Web API)
\end{itemize}

\section{Lớp Giao diện (Presentation Layer -- Cinema.Web)}\label{Sec3.1}

\noindent Lớp giao diện được xây dựng bằng \textbf{ASP.NET Core MVC}. Mỗi chức năng có một Controller tương ứng (Bảng~\ref{table_controller}):

\begin{table}[H]
    \begin{center}
        \begin{tabular}{|l|l|}
            \hline
            \textbf{Controller} & \textbf{Chức năng} \\
            \hline
            \texttt{AccountController}           & Đăng nhập, đăng ký, đăng xuất \\
            \hline
            \texttt{PhimController}              & Danh sách phim, tìm kiếm, chi tiết, chọn ghế \\
            \hline
            \texttt{HoaDonController}            & Xem giỏ hàng (\texttt{GioHang}), xác nhận thanh toán, E-Ticket \\
            \hline
            \texttt{NguoiDungController}         & Hồ sơ (Profile), cập nhật thông tin, đổi mật khẩu \\
            \hline
            \texttt{DichVuController}            & Danh sách dịch vụ bắp nước/combo \\
            \hline
            \texttt{Admin/PhimController}        & Quản trị phim: CRUD + upload ảnh lên \texttt{/images/phim/} \\
            \hline
            \texttt{Admin/SuatChieuController}   & Quản trị suất chiếu: CRUD, tự tính giờ kết thúc \\
            \hline
            \texttt{Admin/DichVuController}      & Quản trị dịch vụ: CRUD (thêm mới trong phiên này) \\
            \hline
            \texttt{Admin/DoanhThuController}    & Báo cáo doanh thu theo phim + biểu đồ Chart.js \\
            \hline
            \texttt{Admin/HomeController}        & Dashboard: tổng phim, vé bán, doanh thu, khách hàng \\
            \hline
        \end{tabular}
    \end{center}
    \caption{Danh sách Controller trong hệ thống}\label{table_controller}
\end{table}

\noindent Lớp này \textbf{không} truy cập cơ sở dữ liệu trực tiếp, mà chỉ gọi qua các interface của tầng BUS được \textbf{Dependency Injection} (DI) tiêm vào constructor thông qua \texttt{Program.cs}. Session được sử dụng để lưu trữ giỏ hàng tạm thời.

\begin{rem}
    Thuộc tính \texttt{[AdminAuthorize]} là một filter tùy chỉnh kiểm tra \texttt{Session["Role"] == "Admin"} trước mỗi request đến khu vực quản trị. Nếu không đủ quyền, hệ thống tự động chuyển hướng về trang đăng nhập.
\end{rem}

\section{Lớp Nghiệp vụ (Business Logic Layer -- Cinema.BUS)}\label{Sec3.2}

\noindent Tầng BUS đóng gói toàn bộ \textbf{logic nghiệp vụ} và được định nghĩa qua các interface (Bảng~\ref{table_bus}):

\begin{table}[H]
    \begin{center}
        \begin{tabular}{|l|l|}
            \hline
            \textbf{Interface / Class} & \textbf{Chức năng chính} \\
            \hline
            \texttt{IPhimBUS / PhimBUS}           & CRUD phim, tìm kiếm, gọi SP, lấy ghế theo suất \\
            \hline
            \texttt{IHoaDonBUS / HoaDonBUS}       & Lưu hóa đơn + vé + dịch vụ trong Transaction \\
            \hline
            \texttt{INguoiDungBUS / NguoiDungBUS} & Đăng ký, đăng nhập, xem profile, đổi mật khẩu \\
            \hline
            \texttt{IDichVuBUS / DichVuBUS}       & Lấy danh sách dịch vụ theo ID \\
            \hline
        \end{tabular}
    \end{center}
    \caption{Các interface và class trong tầng BUS}\label{table_bus}
\end{table}

\section{Đối tượng trao đổi dữ liệu (Cinema.DTO)}\label{Sec3.3}

\noindent Tầng DTO chứa các lớp \textbf{chỉ mang dữ liệu}, không có logic nghiệp vụ, dùng để trao đổi giữa các lớp:

\begin{table}[H]
    \begin{center}
        \begin{tabular}{|l|l|}
            \hline
            \textbf{DTO Class} & \textbf{Mô tả} \\
            \hline
            \texttt{PhimDTO}            & Thông tin phim + danh sách suất chiếu hiển thị \\
            \hline
            \texttt{SuatChieuDTO}       & Mã suất, giờ bắt đầu \\
            \hline
            \texttt{SuatChieuHienThiDTO} & Mã suất, chuỗi giờ hiển thị (dạng \texttt{hh:mm}) \\
            \hline
            \texttt{GheDTO}             & Mã ghế, tên ghế, loại ghế, trạng thái đã đặt \\
            \hline
            \texttt{DichVuDTO}          & Thông tin dịch vụ, số lượng chọn, thành tiền \\
            \hline
            \texttt{DichVuChonDTO}      & Dịch vụ người dùng chọn (MaDV, SoLuong, DonGia) \\
            \hline
            \texttt{CartItemDTO}        & Giỏ hàng: phim, ghế, dịch vụ, tổng tiền \\
            \hline
            \texttt{HoaDonDTO}          & Hóa đơn đầy đủ cho E-Ticket (phim, ghế, dịch vụ) \\
            \hline
            \texttt{NguoiDungDTO}       & Thông tin người dùng + lịch sử hóa đơn \\
            \hline
            \texttt{DoiMatKhauRequest}  & Mật khẩu cũ, mật khẩu mới, xác nhận mật khẩu mới \\
            \hline
        \end{tabular}
    \end{center}
    \caption{Danh sách các lớp DTO}\label{table_dto}
\end{table}

\section{Lớp Truy cập Dữ liệu (Data Access Layer -- Cinema.DAL)}\label{Sec3.4}

\noindent Tầng DAL sử dụng hai phương thức truy cập dữ liệu song song:

\begin{itemize}
    \item \textbf{Entity Framework Core (Database First):} Lớp \texttt{QuanLyRapPhimContext} kế thừa \texttt{DbContext} định nghĩa 10 \texttt{DbSet<T>} và ánh xạ cột qua \texttt{OnModelCreating()}. Chuỗi kết nối đọc từ \texttt{appsettings.json} qua DI.
    \item \textbf{ADO.NET:} Lớp \texttt{CinemaAdoNetDAL} (triển khai \texttt{ICinemaAdoNetDAL}) cung cấp các phương thức truy cập thuần túy bằng \texttt{SqlConnection}, \texttt{SqlCommand}, \texttt{SqlDataAdapter}, \texttt{DataTable} và \texttt{DataSet}.
\end{itemize}

\chapter{Thực thi Kỹ thuật và Công nghệ}

\section{Entity Framework Core: Database First và Code First}\label{Sec4.1}

\subsection{Database First}\label{sub4.1.1}

\noindent Toàn bộ 10 entity trong \texttt{Cinema.DAL/Models/} được scaffold tự động từ SQL Server bằng lệnh \texttt{dotnet ef dbcontext scaffold}. \texttt{QuanLyRapPhimContext} cấu hình toàn bộ ràng buộc, index và quan hệ qua \texttt{OnModelCreating()}.

\subsection{Code First}\label{sub4.1.2}

\noindent Ngoài Database First, nhóm cũng minh họa \textbf{Code First} bằng cách tạo model \texttt{NhatKyHeThong} với Data Annotation và đăng ký vào \texttt{DbContext}:

\begin{lstlisting}[style=csharpstyle, caption={Model Code First -- \texttt{NhatKyHeThong.cs}}]
[Table("NhatKyHeThong")]
public class NhatKyHeThong
{
    [Key]
    [DatabaseGenerated(DatabaseGeneratedOption.Identity)]
    public int MaNhatKy { get; set; }

    [Required]
    [MaxLength(100)]
    public string HanhDong { get; set; } = string.Empty;

    [MaxLength(200)]
    public string? MoTa { get; set; }

    [MaxLength(50)]
    public string? TaiKhoan { get; set; }

    public DateTime ThoiGian { get; set; } = DateTime.Now;
}
\end{lstlisting}

\noindent Để tạo bảng trong database, chạy Migration:
\begin{lstlisting}[style=sqlstyle, caption={Lệnh Migration Code First}]
dotnet ef migrations add AddNhatKyHeThong 
    --project Cinema.DAL 
    --startup-project Cinema.Web
dotnet ef database update 
    --project Cinema.DAL 
    --startup-project Cinema.Web
\end{lstlisting}

\section{LINQ: LINQ to Objects và LINQ to XML}\label{Sec4.2}

\subsection{LINQ to Objects}\label{sub4.2.1}

\begin{thm}\label{thmlinq}
    Toàn bộ thao tác dữ liệu trong tầng BUS sử dụng \textbf{LINQ to Objects} kết hợp \textbf{LINQ to Entities} (qua Entity Framework Core), đảm bảo kiểu an toàn và hạn chế SQL Injection.
\end{thm}

\noindent Ví dụ lấy chi tiết hóa đơn đầy đủ với \textbf{Eager Loading} (\texttt{Include}/\texttt{ThenInclude}):

\begin{lstlisting}[style=csharpstyle, caption={LINQ Eager Loading trong \texttt{HoaDonBUS.LayChiTietHoaDonFull()}}]
var hd = _context.HoaDons
    .Include(h => h.Ves)
        .ThenInclude(v => v.MaGheNavigation)
    .Include(h => h.Ves)
        .ThenInclude(v => v.MaSuatNavigation)
            .ThenInclude(s => s!.MaPhimNavigation)
    .Include(h => h.ChiTietDvs)
        .ThenInclude(ct => ct.MaDvNavigation)
    .FirstOrDefault(h => h.MaHd == maHD);
\end{lstlisting}

\subsection{LINQ to XML}\label{sub4.2.2}

\noindent Hệ thống cung cấp endpoint \texttt{GET /api/phimapi/export-xml} để xuất danh sách phim ra XML, và \texttt{POST /api/phimapi/import-xml} để nhập từ XML, sử dụng \textbf{LINQ to XML} (\texttt{XDocument}, \texttt{XElement}):

\begin{lstlisting}[style=csharpstyle, caption={LINQ to XML -- Export phim ra XML}]
// Lấy dữ liệu bằng LINQ to Entities
var phims = _db.Phims.ToList();

// Xây dựng tài liệu XML bằng LINQ to XML (query syntax)
XDocument xmlDoc = new XDocument(
    new XDeclaration("1.0", "utf-8", "yes"),
    new XElement("DanhSachPhim",
        new XAttribute("NgayXuat", 
                       DateTime.Now.ToString("yyyy-MM-dd")),
        from p in phims
        select new XElement("Phim",
            new XAttribute("MaPhim", p.MaPhim),
            new XElement("TenPhim", p.TenPhim),
            new XElement("ThoiLuong", p.ThoiLuong ?? 0),
            new XElement("NgayKhoiChieu",
                p.NgayKhoiChieu?.ToString("yyyy-MM-dd") ?? "")
        )
    )
);
return Content(xmlDoc.ToString(), "application/xml");
\end{lstlisting}

\begin{lstlisting}[style=csharpstyle, caption={LINQ to XML -- Import phim từ XML}]
XDocument xmlDoc = XDocument.Parse(xmlContent);

// Dùng LINQ to XML (query syntax) để trích xuất dữ liệu
var importedPhims = (from elem in xmlDoc.Descendants("Phim")
                     select new PhimDTO
                     {
                         TenPhim = (string?)elem.Element("TenPhim")
                                   ?? "Khong ten",
                         ThoiLuong = (int?)elem.Element("ThoiLuong") 
                                     ?? 0,
                     }).ToList();
\end{lstlisting}

\section{ADO.NET: Mô hình kết nối và phi kết nối}\label{Sec4.3}

\noindent Lớp \texttt{CinemaAdoNetDAL} (triển khai \texttt{ICinemaAdoNetDAL}) minh họa cả hai mô hình ADO.NET:

\subsection{Mô hình phi kết nối -- DataTable}\label{sub4.3.1}

\begin{lstlisting}[style=csharpstyle, caption={ADO.NET -- SqlDataAdapter + DataTable (phi kết nối)}]
public DataTable LayDanhSachPhimAdoNet()
{
    DataTable dt = new DataTable();
    using (SqlConnection conn = new SqlConnection(_connectionString))
    {
        string query = "SELECT MaPhim, TenPhim, ThoiLuong, " +
                       "GioiHanTuoi, NgayKhoiChieu FROM Phim";
        using (SqlDataAdapter adapter = new SqlDataAdapter(query, conn))
        {
            adapter.Fill(dt); // Phi kết nối: Fill vào bộ nhớ
        }
    }
    return dt; // Endpoint: GET /api/phimapi/adonet
}
\end{lstlisting}

\subsection{Mô hình phi kết nối -- DataSet (nhiều DataTable)}\label{sub4.3.2}

\begin{lstlisting}[style=csharpstyle, caption={ADO.NET -- DataSet chứa Phim + SuatChieu + DataRelation}]
public DataSet LayPhimVaSuatChieu_DataSet()
{
    DataSet ds = new DataSet("CinemaDataSet");
    using (SqlConnection conn = new SqlConnection(_connectionString))
    {
        // DataTable 1: Phim
        using (SqlDataAdapter a1 = new SqlDataAdapter(
                    "SELECT MaPhim, TenPhim, ThoiLuong FROM Phim",
                    conn))
        { a1.Fill(ds, "Phim"); }

        // DataTable 2: SuatChieu
        using (SqlDataAdapter a2 = new SqlDataAdapter(
                    @"SELECT sc.MaSuat, sc.MaPhim, sc.NgayChieu,
                             sc.GioBatDau, sc.GiaVe
                      FROM SuatChieu sc ORDER BY sc.NgayChieu DESC",
                    conn))
        { a2.Fill(ds, "SuatChieu"); }

        // Thiết lập DataRelation giữa 2 bảng trong DataSet
        DataRelation rel = new DataRelation(
            "FK_Phim_SuatChieu",
            ds.Tables["Phim"]!.Columns["MaPhim"]!,
            ds.Tables["SuatChieu"]!.Columns["MaPhim"]!, false);
        ds.Relations.Add(rel);
    }
    return ds; // Endpoint: GET /api/phimapi/adonet-dataset
}
\end{lstlisting}

\subsection{Mô hình kết nối -- SqlDataReader}\label{sub4.3.3}

\begin{lstlisting}[style=csharpstyle, caption={ADO.NET -- SqlDataReader (mô hình kết nối, đọc View)}]
public Dictionary<string, decimal> GetDoanhThuTheoPhimChart()
{
    var result = new Dictionary<string, decimal>();
    using (SqlConnection conn = new SqlConnection(_connectionString))
    {
        string sql = "SELECT TenPhim, TongDoanhThu " +
                     "FROM vw_DoanhThuTheoPhim";
        using (SqlCommand cmd = new SqlCommand(sql, conn))
        {
            conn.Open(); // Kết nối trực tiếp
            using (SqlDataReader reader = cmd.ExecuteReader())
            {
                while (reader.Read())
                {
                    result[reader["TenPhim"].ToString()!] = 
                        Convert.ToDecimal(reader["TongDoanhThu"]);
                }
            }
        }
    }
    return result;
}
\end{lstlisting}

\section{Quản lý Session và Giỏ hàng}\label{Sec4.4}

\noindent Giỏ hàng (\texttt{CartItemDTO}) được serialize thành JSON và lưu vào \textbf{ASP.NET Core Session}. Lớp mở rộng \texttt{SessionExtensions} cung cấp hai phương thức generic:

\begin{lstlisting}[style=csharpstyle, caption={\texttt{SessionExtensions.cs} -- Serialize/Deserialize vào Session}]
public static void Set<T>(this ISession session, string key, T value)
{
    session.SetString(key, JsonSerializer.Serialize(value));
}

public static T Get<T>(this ISession session, string key)
{
    var value = session.GetString(key);
    return value == null 
        ? default 
        : JsonSerializer.Deserialize<T>(value);
}
\end{lstlisting}

\section{Web API RESTful}\label{Sec4.5}

\noindent Hệ thống phát triển hai \texttt{ApiController} tại thư mục \texttt{ApiControllers/}, trả dữ liệu định dạng \textbf{JSON} chuẩn RESTful:

\begin{table}[H]
    \begin{center}
        \begin{tabular}{|c|l|l|}
            \hline
            \textbf{Method} & \textbf{URL} & \textbf{Mô tả} \\
            \hline
            \texttt{GET}    & \texttt{/api/phimapi}              & Lấy phim đang chiếu (gọi SP) \\
            \hline
            \texttt{GET}    & \texttt{/api/phimapi/\{id\}}       & Chi tiết 1 phim \\
            \hline
            \texttt{POST}   & \texttt{/api/phimapi}              & Thêm phim mới, trả \texttt{201 Created} \\
            \hline
            \texttt{PUT}    & \texttt{/api/phimapi/\{id\}}       & Cập nhật phim \\
            \hline
            \texttt{DELETE} & \texttt{/api/phimapi/\{id\}}       & Xóa phim \\
            \hline
            \texttt{GET}    & \texttt{/api/phimapi/adonet}       & Phim qua ADO.NET DataTable \\
            \hline
            \texttt{GET}    & \texttt{/api/phimapi/adonet-dataset} & DataSet (Phim + SuatChieu) \\
            \hline
            \texttt{GET}    & \texttt{/api/phimapi/export-xml}   & Export XML (LINQ to XML) \\
            \hline
            \texttt{POST}   & \texttt{/api/phimapi/import-xml}   & Import phim từ XML \\
            \hline
            \hline
            \texttt{GET}    & \texttt{/api/dichvuapi}            & Lấy tất cả dịch vụ \\
            \hline
            \texttt{GET}    & \texttt{/api/dichvuapi/\{id\}}     & Chi tiết 1 dịch vụ \\
            \hline
            \texttt{POST}   & \texttt{/api/dichvuapi}            & Thêm dịch vụ \\
            \hline
            \texttt{PUT}    & \texttt{/api/dichvuapi/\{id\}}     & Cập nhật dịch vụ \\
            \hline
            \texttt{DELETE} & \texttt{/api/dichvuapi/\{id\}}     & Xóa dịch vụ (kiểm tra FK) \\
            \hline
        \end{tabular}
    \end{center}
    \caption{Danh sách các endpoint Web API}\label{table_api}
\end{table}

\section{Luồng đặt vé và thanh toán}\label{Sec4.6}

\noindent Luồng xử lý hoàn chỉnh từ khi chọn ghế đến khi nhận E-Ticket (Bảng~\ref{table_flow}):

\begin{table}[H]
    \begin{center}
        \begin{tabular}{|c|l|l|}
            \hline
            \textbf{Bước} & \textbf{Hành động} & \textbf{Endpoint} \\
            \hline
            1 & Chọn phim và suất chiếu    & \texttt{GET /Phim/Index} \\
            \hline
            2 & Xem chi tiết, chọn suất   & \texttt{GET /Phim/Details/\{id\}} \\
            \hline
            3 & Chọn ghế, thêm dịch vụ    & \texttt{GET /Phim/ChonGhe?maSuat=\{id\}} \\
            \hline
            4 & Lưu giỏ hàng vào Session  & \texttt{POST /Phim/LuuGheVaoSession} (AJAX) \\
            \hline
            5 & Xem giỏ hàng              & \texttt{GET /HoaDon/GioHang} \\
            \hline
            6 & Xác nhận thanh toán       & \texttt{POST /HoaDon/XacNhanThanhToan} (AJAX) \\
            \hline
            7 & Nhận E-Ticket + QR Code   & \texttt{GET /HoaDon/ThanhToanThanhCong?id=\{maHD\}} \\
            \hline
        \end{tabular}
    \end{center}
    \caption{Luồng đặt vé và thanh toán}\label{table_flow}
\end{table}

% ============================================================
\chapter{Kết luận và Hướng phát triển}

\noindent Đồ án đã hoàn thành các mục tiêu đề ra:

\begin{itemize}
    \item \textbf{Lập trình CSDL T-SQL:} Xây dựng 10 bảng quan hệ trong SQL Server với đầy đủ View (\texttt{vw\_DoanhThuTheoPhim}), Function (\texttt{fn\_TinhTongTienHoaDon}), Stored Procedure (\texttt{sp\_LayDanhSachPhimDangChieu}), Trigger (\texttt{trg\_NganXoaHoaDonCoVe}) và Transaction (EF Core \texttt{BeginTransaction/Commit/Rollback}).
    \item \textbf{Kiến trúc 3 lớp:} Tách biệt hoàn toàn Presentation (Cinema.Web), Business Logic (Cinema.BUS) và Data Access (Cinema.DAL), giao tiếp qua DTO và Interface.
    \item \textbf{ADO.NET:} Triển khai cả mô hình kết nối (\texttt{SqlDataReader}) và phi kết nối (\texttt{SqlDataAdapter}, \texttt{DataTable}, \texttt{DataSet} với \texttt{DataRelation}).
    \item \textbf{LINQ:} Sử dụng LINQ to Objects/Entities xuyên suốt tầng BUS và LINQ to XML cho tính năng Export/Import dữ liệu.
    \item \textbf{Entity Framework:} Database First (scaffold 10 entity) và Code First (model \texttt{NhatKyHeThong} với Data Annotation + Migration).
    \item \textbf{Web API RESTful:} 2 Controller API (\texttt{PhimApiController}, \texttt{DichVuApiController}) với 14 endpoint trả JSON chuẩn REST (GET/POST/PUT/DELETE).
    \item \textbf{Tính năng người dùng:} Đặt vé, chọn ghế, dịch vụ combo, thanh toán, E-Ticket QR, hồ sơ cá nhân và \textbf{đổi mật khẩu}.
    \item \textbf{Khu vực Admin:} CRUD phim, suất chiếu, dịch vụ, báo cáo doanh thu theo phim.
\end{itemize}

\noindent\textbf{Hướng phát triển tiếp theo:}

\begin{itemize}
    \item Tích hợp \textbf{cổng thanh toán trực tuyến} (VNPay, MOMO) thay vì thanh toán giả lập.
    \item Áp dụng \textbf{JWT Authentication} thay thế Session để hỗ trợ REST API cho ứng dụng mobile.
    \item Thêm module \textbf{quản lý phòng chiếu và ghế động} (tự động sinh ghế theo sức chứa phòng).
    \item Áp dụng \textbf{mã hóa mật khẩu} (BCrypt/PBKDF2) thay vì lưu plaintext hiện tại.
    \item Triển khai \textbf{CI/CD} trên Azure / AWS và logging tập trung với Serilog.
    \item Viết \textbf{Unit Test} cho tầng BUS sử dụng xUnit và Moq.
\end{itemize}

% ============================================================
\begin{thebibliography}{99}
\addcontentsline{toc}{chapter}{TÀI LIỆU THAM KHẢO}

\bibitem{ardalis2023}
Steve ``Ardalis'' Smith, \textit{Architecting Modern Web Applications with ASP.NET Core and Azure}, Microsoft, 2023. Available: \texttt{https://dotnet.microsoft.com/learn/aspnet/architecture}

\bibitem{peres2018}
Ricardo Peres, \textit{Entity Framework Core Succinctly}, Syncfusion, 2018.

\bibitem{vieira2012}
Robert Vieira, \textit{Beginning Microsoft SQL Server 2012 Programming}, Wrox, 2012.

\bibitem{msefcore}
Microsoft, \textit{Entity Framework Core Documentation}. Available: \texttt{https://learn.microsoft.com/en-us/ef/core/}

\bibitem{msaspnet}
Microsoft, \textit{ASP.NET Core MVC Documentation}. Available: \texttt{https://learn.microsoft.com/en-us/aspnet/core/mvc/}

\bibitem{msadonet}
Microsoft, \textit{ADO.NET Overview -- DataSet, DataAdapter}. Available: \texttt{https://learn.microsoft.com/en-us/dotnet/framework/data/adonet/}

\bibitem{mslinq}
Microsoft, \textit{LINQ to XML Overview (C\#)}. Available: \texttt{https://learn.microsoft.com/en-us/dotnet/standard/linq/linq-xml-overview}

\end{thebibliography}

\end{document}
